% Limitations: one place, one sentence each (REWRITE_PLAN_2026-09-26 section 5). Not repeated in Results.

\subsection{Limitations}

Each demographic cell has one persona, so a persona's own quirks cannot be separated from its cell.
Transgender and multiracial personas have no NHANES comparison group, so they are not tested at levels 2 and 3.
At NHANES precision level 2 can detect distortion and cannot certify faithfulness.
No clinician reviewed individual answer vectors, and that review is the next test of level 1.
Of the clinical rows, 1,575 are not the December run's own output, 1,345 of them came from a prompt that cannot be established, and no per-model verdict changes without them. % R: corpus_v3_provenance.csv; 87_row_sources.csv
The reissue line added to refused calls was not logged per attempt, so the rows produced under it cannot be identified.
The personas carry no age, so the simulated population's age mix is unknown and is not matched to the reference.
The PHQ-8 items were presented in one fixed order.
Each attribute level was given to the model in one wording, so the verdicts apply to these wordings.
No analysis was registered with a public registry.

\subsection{Conclusion}

This paper proposes the validity ladder as a protocol that other researchers can apply to any simulated population and extend to new instruments.
In the worked example, four models generated 28,800 PHQ-8 assessments, and every model failed every level above the gate. % R: corpus_v3_provenance.csv
Pseudo-models built from survey respondents passed levels 1 and 3 and the level-4 structure rules, and planted failures at those levels were detected.
Level 2 detected planted distortion but cannot certify a faithful model at survey precision.
A researcher who uses language models to represent a population should therefore run each rung before reporting a claim that the rung supports.
The code and a receipt for every reported number are public, so the ladder can be run on a new corpus without rebuilding it.
